\chapter{Initramfs}
\label{ch:initramfs}

The initial RAM filesystem (initramfs) is a compressed cpio archive
embedded in the kernel binary that serves as the root filesystem during
and after boot. Unlike traditional initrd (which uses a block device and
a filesystem driver), initramfs uses tmpfs and is managed by the kernel's
VFS layer directly.

\section{What is Initramfs?}

Initramfs is the kernel's mechanism for providing an initial root
filesystem before the real root filesystem is mounted. In Veilbox,
there is no separate real root — the initramfs \emph{is} the root
filesystem, permanently.

Key characteristics:
\begin{itemize}[nosep]
  \item \textbf{Memory-based}: The root filesystem lives entirely in RAM
        (tmpfs). There is no block device backing it.
  \item \textbf{Compressed}: The initramfs archive is gzip-compressed
        (or xz-compressed, depending on kernel config) inside the kernel
        binary. The kernel decompresses it at boot.
  \item \textbf{Embedded}: The archive is linked into the kernel image
        during the build, not loaded separately. The kernel and its
        rootfs are a single file.
  \item \textbf{Read-only at runtime}: Because the initramfs is embedded
        in the kernel binary, modifications to the root filesystem are
        lost on reboot. Persistent data is stored on a separate state
        partition.
\end{itemize}

\section{The CPIO Archive Format}

CPIO (Copy In/Output) is an archival format like tar but simpler. Each
file in the archive is preceded by a 110-byte header:

\begin{lstlisting}
struct cpio_header {
    char magic[6];      // "070701" for new ASCII format
    char ino[8];        // Inode number (hex)
    char mode[8];       // File mode and type (hex)
    char uid[8];        // Owner UID (hex)
    char gid[8];        // Owner GID (hex)
    char nlink[8];      // Number of hard links (hex)
    char mtime[8];      // Modification time (hex)
    char filesize[8];   // Size of file data (hex)
    char devmajor[8];   // Major device number (hex)
    char devminor[8];   // Minor device number (hex)
    char rdevmajor[8];  // Major for device nodes (hex)
    char rdevminor[8];  // Minor for device nodes (hex)
    char namesize[8];   // Pathname length including null (hex)
    char check[8];      // Checksum (set to "00000000" if not used)
};
\end{lstlisting}

\section{The Kernel File List Format}

Instead of building a cpio archive manually, Veilbox uses the kernel's
internal file list format. This is a text file that the kernel's build
system (\texttt{scripts/gen\_initramfs\_list.sh}) converts to a cpio
archive automatically.

\subsection{File List Syntax}

Each line in the file list specifies one filesystem entry:

\begin{lstlisting}
# Comments start with #
dir   /proc          0755 0 0
dir   /sys           0755 0 0
dir   /dev           0755 0 0
file  /bin/busybox   0755 0 0 path/to/busybox
slink /bin/sh        /bin/busybox    0755 0 0
nod   /dev/console   0600 0 0 c 5 1
nod   /dev/null      0666 0 0 c 1 3
\end{lstlisting}

The five entry types are:

\begin{longtable}{|p{2cm}|p{4cm}|p{6cm}|}
\hline
\textbf{Type} & \textbf{Syntax} & \textbf{Description} \\
\hline
dir & \texttt{dir path mode uid gid} & Creates a directory. \\
nod & \texttt{nod path mode uid gid type major minor} & Creates a device node (c=char, b=block). \\
slink & \texttt{slink path target mode uid gid} & Creates a symbolic link. \\
pipe & \texttt{pipe path mode uid gid} & Creates a FIFO (named pipe). \\
sock & \texttt{sock path mode uid gid} & Creates a Unix domain socket. \\
file & \texttt{file path mode uid gid source} & Copies a file from the host filesystem. \\
\hline
\end{longtable}

\subsection{Generating the File List}

The build script uses the kernel's provided script:

\begin{lstlisting}[language=bash]
gen_initramfs_list() {
    scripts/gen_initramfs_list.sh rootfs/ > initramfs.list

    # Add essential device nodes not created by rootfs/
    echo "nod /dev/console 0600 0 0 c 5 1"   >> initramfs.list
    echo "nod /dev/tty0    0600 0 0 c 4 0"   >> initramfs.list
    echo "nod /dev/ttyS0   0600 0 0 c 4 64"  >> initramfs.list
    echo "nod /dev/null    0666 0 0 c 1 3"   >> initramfs.list
    echo "nod /dev/zero    0666 0 0 c 1 5"   >> initramfs.list

    # Sort for deterministic ordering
    sort -o initramfs.list initramfs.list
}
\end{lstlisting}

\section{How the Kernel Processes Initramfs}

During boot, the kernel performs these steps:

\begin{enumerate}[nosep]
  \item \textbf{Decompression}: After the kernel decompresses itself and
        sets up basic memory management, it locates the embedded cpio
        archive (identified by the magic bytes \texttt{070701} at a known
        offset in the init section).
  \item \textbf{Extraction}: The kernel calls \texttt{populate\_rootfs()}
        in \texttt{init/initramfs.c}. This function reads the cpio archive
        and creates the corresponding directory entries, files, symlinks,
        and device nodes in a tmpfs filesystem.
  \item \textbf{Mounting}: The tmpfs root is mounted at \texttt{/}. At
        this point, the kernel has a fully functional root filesystem in
        memory.
  \item \textbf{Executing init}: The kernel locates the init binary
        (\texttt{/init} or \texttt{/sbin/init}) and executes it as PID 1.
        In Veilbox, this is BusyBox init.
\end{enumerate}

\section{Initramfs Size Considerations}

The initramfs for Veilbox is approximately 90 MB uncompressed and 30 MB
compressed. This is large for an initramfs but necessary because it
contains the complete root filesystem including containerd, runc, and
nerdctl binaries.

\begin{longtable}{|p{4cm}|p{3cm}|p{3cm}|p{3cm}|}
\hline
\textbf{Component} & \textbf{Size (raw)} & \textbf{Size (compressed)} & \textbf{Notes} \\
\hline
BusyBox binary & 1.2 MB & 0.5 MB & Static binary \\
containerd & 35 MB & 12 MB & Dynamic binary \\
runc & 10 MB & 3.5 MB & Static binary \\
nerdctl & 25 MB & 8 MB & Dynamic binary \\
Dropbear & 0.5 MB & 0.2 MB & Static binary \\
Shared libraries & 15 MB & 5 MB & glibc, libseccomp \\
Config files & 50 KB & 10 KB & inittab, rcS, etc. \\
Device nodes & 0 KB & 0 KB & Created by kernel \\
\hline
\textbf{Total} & \textbf{~87 MB} & \textbf{~29 MB} & \\
\hline
\end{longtable}

\section{Device Nodes}

Device nodes in the initramfs are specified in the file list, not created
by a running udev or mdev daemon. The following device nodes are essential:

\begin{longtable}{|p{3cm}|p{2cm}|p{2cm}|p{2cm}|p{4cm}|}
\hline
\textbf{Device} & \textbf{Major} & \textbf{Minor} & \textbf{Type} & \textbf{Purpose} \\
\hline
/dev/console & 5 & 1 & char & Kernel console \\
/dev/tty0 & 4 & 0 & char & VGA text console \\
/dev/ttyS0 & 4 & 64 & char & Serial console \\
/dev/null & 1 & 3 & char & Data sink \\
/dev/zero & 1 & 5 & char & Zero-filled reads \\
/dev/random & 1 & 8 & char & Random number gen \\
/dev/urandom & 1 & 9 & char & Non-blocking random \\
/dev/kmsg & 1 & 11 & char & Kernel log access \\
/dev/ptmx & 5 & 2 & char & PTY multiplexer \\
\hline
\end{longtable}


\section{Device Nodes Explained}

\subsection{Major and Minor Numbers}

Every device node in Linux has a major number (identifying the device driver)
and a minor number (identifying the specific device instance):

\begin{lstlisting}
# Character devices
/dev/console  -> major 5, minor 1  (console driver)
/dev/tty0     -> major 4, minor 0  (TTY driver, first VGA console)
/dev/ttyS0    -> major 4, minor 64 (TTY driver, first serial port)
/dev/null     -> major 1, minor 3  (memory device, null sink)
/dev/zero     -> major 1, minor 5  (memory device, zero source)
/dev/random   -> major 1, minor 8  (kernel random number generator)
/dev/kmsg     -> major 1, minor 11 (kernel log messages)
\end{lstlisting}

The kernel creates device nodes from the initramfs file list entries.
Each \texttt{nod} line in the list specifies:

\begin{lstlisting}
nod <path> <permissions> <uid> <gid> <type> <major> <minor>
\end{lstlisting}

Where \texttt{type} is \texttt{c} for character devices or \texttt{b} for
block devices.

\subsection{devtmpfs vs. Static Nodes}

Veilbox uses a hybrid approach:

\begin{enumerate}[nosep]
  \item \textbf{Static nodes in initramfs}: The initramfs file list
        pre-creates critical device nodes (\texttt{/dev/console},
        \texttt{/dev/null}, \texttt{/dev/tty0}, \texttt{/dev/ttyS0}).
        These are needed before udev or devtmpfs starts.
  \item \textbf{devtmpfs at runtime}: The \texttt{mount -t devtmpfs}
        command in rcS mounts a device-managed tmpfs that automatically
        creates nodes for all kernel-registered devices.
\end{enumerate}

This ensures that essential devices exist even before devtmpfs is mounted,
while devtmpfs provides complete coverage for all kernel devices.

\section{Initramfs Compression Options}

The kernel supports several compression algorithms for the embedded cpio
archive:

\begin{longtable}{|p{3cm}|p{3cm}|p{3cm}|p{3cm}|}
\hline
\textbf{Algorithm} & \textbf{Compression Ratio} & \textbf{Speed} & \textbf{Kernel Option} \\
\hline
gzip & 3:1 & Fast & CONFIG\_RD\_GZIP \\
bzip2 & 4:1 & Slow & CONFIG\_RD\_BZIP2 \\
lzma & 5:1 & Very slow & CONFIG\_RD\_LZMA \\
xz & 5:1 & Very slow & CONFIG\_RD\_XZ \\
lzo & 2.5:1 & Very fast & CONFIG\_RD\_LZO \\
lz4 & 2:1 & Extremely fast & CONFIG\_RD\_LZ4 \\
zstd & 3.5:1 & Fast & CONFIG\_RD\_ZSTD \\
\hline
\end{longtable}

Veilbox uses gzip compression (the default) for its balance of speed and
compression ratio. The initramfs compresses from ~90 MB to ~30 MB.
